\documentclass[11pt,a4paper]{article}
\usepackage[margin=27mm]{geometry}
\usepackage{amsmath,amssymb,amsthm,booktabs,longtable,array}
\usepackage[T1]{fontenc}
\usepackage{lmodern,microtype}
\usepackage[colorlinks=true,linkcolor=blue,citecolor=blue,urlcolor=blue]{hyperref}
\newtheorem{proposition}{Proposition}
\newcommand{\E}{\mathbb E}
\newcommand{\Prob}{\mathbb P}
\newcommand{\ind}{\mathbf 1}
\title{Financial statements do not determine a company:  learning, compression and financing in scenario simulation}
\author{AccDSL research supplement}
\date{26 October 2026\\Provisional technical synthesis; source and implementation limits stated locally}
\begin{document}
\maketitle
\begin{abstract}
Financial statements constrain the history of a company without identifying all
its transactions, contractual exposures or future decisions. A simulator must
preserve those distinctions, including differences that later information may
never resolve. The computational problem is therefore to share representations
and calculations while retaining the information needed for future decisions.
This supplement develops conditions for exact sharing, explains why approximate
sharing needs separate policy and threshold arguments, and distinguishes hard
evidence from probabilistic learning. It then connects debt proceeds, promised
payments, recovery and shareholder dilution to dated financing choices. Primary
work on accounting representations, adapted transport, sequential inference,
disclosure and corporate finance supplies both useful methods and precise
reasons that some tempting transfers fail. Exact synthetic cases and inspected
author code support the local conclusions. They do not establish empirical
forecast validity or completion of the general company simulator.
\end{abstract}
\tableofcontents

\section{The company behind the reported numbers}
A company can report the same debt balance under a fixed-rate loan and a
floating-rate loan. It can even have recorded the same historical interest
payments if the floating rate has not yet reset. A future rate shock separates
their cash flows. Two different sequences of purchases and disposals can also
produce the same closing asset balance while leaving different assets available
for production or sale. The first example conceals contractual state beyond
the journal; the second conceals transactions within an accounting aggregate.
Neither ambiguity is removed by constructing one balanced journal.

For a fixed linear reporting map, write reported amounts as $y=Az$, where $z$
contains transaction or account movements. If $Av=0$, then $z$ and $z+v$ have
the same reports. They are both admissible only if they also satisfy the
nonnegativity restrictions, dates, notes, contractual conditions and accounting
rules in the declared model. The null space alone is not a generator of
economically possible companies. More generally, let $h$ contain opening state,
dated events, outstanding contracts, recognition and measurement choices, and
source support. With reporting map $R$ and observed evidence $e$, define
\begin{equation}
\mathcal H(e)=\{h:R(h)\text{ agrees with }e
\text{ and }h\text{ satisfies the declared economic restrictions}\}.
\end{equation}
Agreement includes units, rounding, dates and reporting perimeter. An assertion
that a loan is floating rate is as relevant as a reported numeric total. A
missing note is recorded as missing evidence; its contents are not inferred
from the convenience of a candidate reconstruction.

McCarthy's REA model provides a useful architectural antecedent. Resources,
events and agents, linked through stock-flow, duality and control relations,
organise economic activity before particular accounting views are produced
\cite[pp.~560--570]{mccarthy}. Double entry can be one external view. Claims
can be derived from imbalances or retained as base records when their uses
warrant it. The paper also recognises the practical need for aggregation and
the fact that measurement and periodicity are not settled by a database schema
\cite[pp.~568--576]{mccarthy}. These distinctions support an event-and-contract
representation with journal views. They do not imply that disclosures identify
a unique event history, or supply the likelihood of observing a disclosure.

The proposed history family is therefore a family of hypothetical explanations,
not multiple legal books asserted to belong to the issuer. Revisions to actual
reports form a separate version history. Future behaviour is another separate
choice: even complete historical contracts do not identify how management will
respond to a new liquidity shock. A run binds a history, a behavioural model,
an economic scenario and accounting rules, and retains that combination in its
results. Without defensible weights over histories, the result is a family of
conditional laws or feasible outcomes. Counting equivalent encodings of one
history must not make that history more probable.

\section{Possibilities do not need individual tree nodes}
Suppose twenty contractual clauses are unknown binary facts. The number of
assignments is $2^{20}$, but twenty variables and their joint restrictions need
not require a million copies of the company. Common events can be stored once;
uncertain rates can remain symbolic; alternative clauses can be guards on a
shared calculation. This saves representation and repeated computation. It
does not guarantee that every query is easy: some queries genuinely require
exploring many combinations, and continuous uncertainty is not a finite list.

A tree is useful when showing what is observed and decided over time. A stored
node may represent a region of feasible histories, with outgoing edges denoting
observations or economically meaningful events. A graph can share common
prefixes and equivalent continuations. Distinct provenance is retained even
when execution is shared. The meaning of an edge matters: logical dependence,
source support and stochastic transition are different relations. Drawing the
same graph for them does not establish conditional independence.

Consider a loan whose unknown fixed coupon is either 5 or 7 in each of two
periods. The admissible payment paths are $(5,5)$ and $(7,7)$. Taking the
Cartesian product of the marginal possibilities $\{5,7\}$ at each date adds
$(5,7)$ and $(7,5)$. Those extra paths are wrong for this contract. A change of
terms requires a modification or refinancing event. The same error arises if a
robust dynamic program allows a new adversarial choice of an otherwise fixed
contract each period. A rectangular family can be appropriate for a different
model, but it changes this one. Dynamic consistency of preferences and factual
persistence of a contract must be assessed separately; the learning-under-
ambiguity literature makes updating and restrictions on priors substantive
assumptions \cite{epstein}.

\subsection{Exact sharing must preserve information as well as cash}
Assume a fixed finite horizon and a transition model that is Markov in the
augmented economic state, including persistent terms and time. Let $s$ be
such a finite state and $a$ an action available to a specified
actor. Write $K_s^a(B,o,z)$ for the joint probability that the next state lies
in block $B$, the actor observes $o$, and the requested current output is $z$.
Two states can be placed in the same block when they have the same feasible
actions, the same already available information, the same current outputs,
and the same $K_s^a$ for every admitted action. Any terminal output must also
be constant on each terminal block. Persistent unweighted alternatives are
kept as separate types, with this comparison performed within the relevant
conditional model rather than by inventing a distribution over types.

For a common policy that depends only on the preserved observation history,
these conditions preserve the joint law of future observations and outputs.
To see why, begin at the last date: the terminal output agrees within each
block. One step earlier, the common policy chooses the same action distribution
at indistinguishable histories, and the joint transition into successor block,
observation and output agrees. Summing the common successor law gives the same
continuation law. Backward induction repeats that argument to the starting
date. This is a finite-state argument under the stated joint-kernel condition;
it is not a general theorem that equal balance sheets imply equivalent states.

The observation-quotient work already present in AccDSL supplies exact finite
path comparisons and deliberately rejected partitions for this restriction.
The present supplement reuses that work. A query that adds another output,
policy, horizon or observation can invalidate a previously justified merge.
An analyst may not know the loan type even though management does; policies
then use management's information. Conversely, sampling next year's shock
inside the simulator must not reveal it to management today.


The distinction between a process and its information becomes especially clear
in the work of Beiglb\"ock, Pfl\"ugl and Schrott. Their information process
recursively augments today's value by the conditional law of tomorrow's
information process. Equal laws of suitable Markov lifts characterize their
notion of equivalence of filtered processes
\cite[Definitions~2.6--2.8, Proposition~3.6 and Theorem~3.9]{probview}.
A lift may be large or measure-valued. Its existence does not imply that a
small sufficient state can be recovered from financial statements. Moreover,
transferring strategies requires attention to independent randomization and
the available stochastic basis \cite[Theorem~5.4 and Remark~5.7]{probview}.
An implementer must decide whether an actor may randomize and which observations
are available, rather than deducing these permissions from equal outcome laws.

\subsection{A small example exposes the cost of lost information}
Let tomorrow's payoff driver $Z$ be a fair binary variable. Today an actor
chooses $a\in\{0,1\}$ and receives $\ind\{a=Z\}$. In one model, today's
signal reveals $Z$; in another, the signal is constant. Both models have the
same distribution of $Z$. The optimal expected payoff is 1 in the revealing
model and $1/2$ in the uninformative model. A distance that compares only the
terminal distribution reports no discrepancy while the decision problem
changes substantially. The exact diagnostic enumerates every feasible policy
in both models; it does not allow a separate action for a hidden state.

Nested and adapted transport distances address the temporal issue by
restricting couplings so that matching paths cannot depend on information
that is unavailable at the relevant date. Pflug and Pichler's stability
argument uses convex decision sets, convex loss in decisions and uniform
regularity in uncertainty \cite[Definition~5.1, Theorem~6.1 and proof]{pflug}.
A projected action is formed by conditional averaging. If the only permitted
issuance sizes are 0 and 1, an average of $1/2$ is infeasible. The theorem
therefore does not automatically justify merging discrete issuance choices.

Backhoff-Veraguas and coauthors obtain related stability results for financial
optimization with bounded predictable strategies and regular payoffs
\cite[Theorem~1.5; Lemma~3.7; Theorem~3.10]{adapted}. Their weak transfer
permits a different strategy in the other model. Reusing the same strategy
needs additional regularity. Their semimartingale distance also has a
different cost functional from a simple sum of discrete path deviations; the
paper explains why time refinement matters. Its examples of discontinuous
strategies and exact superhedging delimit the guarantee. These are useful
reasons to formulate the company decision problem precisely, rather than
reasons to attach one distance to every accounting query.

\section{Approximation needs the error relevant to the decision}
\subsection{Thresholds require control near the boundary}
Default, covenant breach and a margin call are threshold events. Let $X$ be
exact cash at a decision date, $Y$ its approximation under a specified joint
coupling, and $c$ the contractual threshold. All three are in the same currency
and date basis. A cash approximation of $-\epsilon$ by $+\epsilon$ has error
$2\epsilon$, yet changes $\ind\{X<0\}$ from one to zero. Small cash error
alone does not imply small default-probability error.

\begin{proposition}
If $\E|X-Y|<\infty$, then for every $\delta>0$,
\begin{equation}\label{eq:threshold}
|\Prob(X<c)-\Prob(Y<c)|
\leq \Prob(|X-c|\leq\delta)+\frac{\E|X-Y|}{\delta}.
\end{equation}
\end{proposition}
\begin{proof}
The absolute difference in event probabilities is at most the probability
that their indicators disagree. If $|X-c|>\delta$ and $|X-Y|<\delta$, $X$
and $Y$ lie on the same side of $c$. Consequently disagreement is contained
in $\{|X-c|\leq\delta\}\cup\{|X-Y|\geq\delta\}$. The union bound
and Markov's inequality give (\ref{eq:threshold}).
\end{proof}
The first term is probability mass close to the covenant; the second is the
probability allowed for a sufficiently large approximation error. This local
derivation is independent of any claim that a particular algorithm supplies
the required coupling error. Write $e=\E|X-Y|$. If $e=0$, then $X=Y$ almost surely and the threshold probabilities agree.
For $e>0$, if $X$ has a density bounded by $M>0$ near $c$,
the first term is at most $2M\delta$. Choosing
$\delta=\sqrt{e/(2M)}$, when that choice lies within the neighborhood where
the density bound holds, gives $2\sqrt{2Me}$. Atoms at the threshold and
uncontrolled conditional distributions prevent this simplification.

For several dates, a union bound can bound the probability that any breach
indicator differs by the sum of the per-date bounds, provided the relevant
coupling and cash errors are supplied. If a different breach decision changes
subsequent actions and cash flows, those cash errors must already account for
that feedback; an open-loop cash comparison does not establish them. When
there are no weights over histories, the appropriate output is instead an
uncertain boundary region or certified feasible-set containment. A probability
bound cannot be manufactured from the proportion of enumerated witnesses.

\subsection{Close optimal values do not validate a transferred policy}
Let $J(\pi)$ and $\widetilde J(\pi)$ be expected costs in two models, with
the same admissible policy class $\Pi$, finite expected costs and finite optimal values. This comparison presupposes that a
policy can be implemented with the information and feasible actions in both
models. Suppose a uniform bound
$\sup_{\pi\in\Pi}|J(\pi)-\widetilde J(\pi)|\leq\varepsilon$ holds.
If $\widehat\pi$ is $\eta$-optimal for the approximate problem, then
\begin{align}
J(\widehat\pi)&\leq\widetilde J(\widehat\pi)+\varepsilon\nonumber\\
&\leq\inf_{\pi\in\Pi}\widetilde J(\pi)+\eta+\varepsilon\nonumber\\
&\leq\inf_{\pi\in\Pi}J(\pi)+\eta+2\varepsilon.
\end{align}
The last step follows by applying the uniform bound to an arbitrarily
near-optimal policy of the exact problem, so attainment is unnecessary.
The result is a regret bound for a common feasible policy. If one only knows
that the two infima are close, the first inequality is missing. If a policy
lift changes the policy class, a bound for that lift and the relevant comparison
class is needed. This is why reporting a small optimized residual or close
portfolio value alone does not establish that a compressed financing policy
will work in the original company.

\subsection{Even sampling can accidentally reveal the future}
With continuous first-period observations, independent sampled paths typically
have distinct first values. If each stored first value has only its single
sampled continuation, that value becomes an identifier for the sampled future.
An optimizer can exploit this unintended information. More paths do not by
themselves repair the information structure of the empirical measure.
Backhoff and coauthors give an explicit optimal-stopping counterexample
\cite[Example~7.1]{estimation}. Their adapted empirical measure quantizes
observations so that paths share histories. Strong consistency applies to
iid complete paths on a fixed finite horizon and bounded state space;
quantitative rates additionally require Lipschitz conditional kernels
\cite[Definitions~1.1--1.2, Theorems~1.3, 1.5, 1.7]{estimation}.

A Markov modification improves the horizon exponent only when the process
is in fact Markov \cite[Theorem~6.1]{estimation}; constants still depend on
the horizon. A company's current balance sheet generally is not a sufficient
Markov state when hidden debt terms or asset vintages matter. Nor are reports
from one issuer an iid sample of complete company paths. The method suggests
how a deliberately generated simulation sample can retain information. Its
rates cannot simply be assigned to an observed issuer panel.

\section{Learning without erasing persistent possibilities}
Let $\theta$ denote fixed but unknown contract terms and $X_t$ the evolving
economic state. Their roles differ. The transition for $X_t$ can depend on
$\theta$, but absence of evidence does not authorize resampling $\theta$
as a new physical contract. Management, lenders and the analyst can have
different observation histories $I_t^a$; each action must be measurable with
respect to the information of the actor choosing it. The simulator keeps
physical state separate from what each actor believes about that state.

For hard evidence $e_{t+1}$ that concerns the same uncorrected historical
assertions, update the feasible family by intersection:
\begin{equation}
\mathcal H_{t+1}=\{h\in\mathcal H_t:h\text{ agrees with }e_{t+1}\}.
\end{equation}
Empty intersection exposes conflict, model misspecification or invalid
evidence. It is not a request to replace the result with the nearest convenient
history. A correction is different: rebuild the evidence view after
superseding the erroneous assertion, preserve the old view, and recompute.
Conjoining both versions can create a false contradiction.

For a probabilistic model, first propagate the joint law of state and fixed
terms under the declared action and transition model. If $\mu^-_{t+1}$ is
that predicted law and $L(o\mid x,\theta,I_t)$ is the observation likelihood,
the update is
\begin{equation}\label{eq:bayes}
\mu_{t+1}(dx,d\theta)=
\frac{L(o_{t+1}\mid x,\theta,I_t)\,\mu^-_{t+1}(dx,d\theta)}
{\int L(o_{t+1}\mid u,\vartheta,I_t)\,\mu^-_{t+1}(du,d\vartheta)}.
\end{equation}
The denominator must be positive and finite. Densities for exact continuous
observations require the corresponding conditional model; treating a
zero-probability point as an ordinary positive-probability event is invalid.
If the likelihood is constant across admissible types, posterior odds are
unchanged. An uninformative sequence need never resolve the uncertainty.

A set of priors or likelihoods produces a set of updates by applying
(\ref{eq:bayes}) to each admissible joint model. Coupled parameter and
likelihood restrictions remain coupled. For a binary type with prior
$p\in[1/4,3/4]$ and likelihoods $4/5$ and $1/5$, the posterior is
$4p/(1+3p)$. Its derivative $4/(1+3p)^2$ is positive, so the posterior range
is $[4/7,12/13]$. These are derived bounds for this specified likelihood,
not confidence bounds inferred from financial reports. With no likelihood
justified, hard-compatible families and sensitivity analysis remain available.

\subsection{A recursive worst case can describe a different company}
Bayesian conditioning and a fresh worst-case choice at every date answer
different questions. Suppose a persistent binary type $\theta$ is either zero
or one, and tomorrow a fair binary signal $X$ unrelated to that type is observed.
The admitted laws are $Q_0$ and $Q_1$: under $Q_i$, $\theta=i$ with certainty
and $X$ remains fair. For the loss $Z=\ind\{X=\theta\}$, either law gives
expected loss $1/2$, so
\begin{equation}
\max_{Q\in\{Q_0,Q_1\}}\E_Q[Z]=\tfrac12,\qquad
\E_X\!\left[\max_{i\in\{0,1\}}\E_{Q_i}[Z\mid X]\right]=1.
\end{equation}
The second expression selects type zero after signal zero and type one after
signal one. It therefore evaluates a new law with $\theta=X$. That law is
absent even from the convex hull of the original two laws: every mixture
retains independence between the signal and type. No type has physically
changed along a realized path, yet dependence has changed across possible
paths. Preserving path support alone would miss the error.

Shapiro defines rectangularity through equality between static worst-case
evaluation and its nested conditional counterpart
\cite[Definition~2.1, Theorem~2.1 and proof, Section~5]{shapiro}.
His finite construction pastes conditional distributions across histories; it
can enlarge the family. A dynamic program for that enlarged family solves
a changed robustness question. Its convenience cannot justify claiming that
it solved the original persistent-type problem.

Terminology also requires care. Pichler and Shapiro call a product family a
rectangular setting, but explicitly show that its original worst-case
functional and its composite conditional functional can differ
\cite[Section~4.2, Example~4.7]{pichlershapiro}. Computing a conditional
law-based risk measure can also differ from conditioning an unconditional
ambiguity set. Take six equiprobable states and an observed three-state
group with losses $0,1,2$. Conditional expected shortfall at level $1/2$
within that group is $5/3$. The unconditional six-state risk envelope permits
probabilities at most $1/3$ per state. It can assign $1/6$ to the loss-2 state,
zero to the other two group members, and $5/18$ to each state outside the
group. This is an admissible probability law; after conditioning on the group
it assigns probability one to loss 2. The conditional worst case over that
envelope is therefore 2. The difference is between two precise conditional
risk definitions, not an arithmetic error
\cite[Example~3.2 and equations~(3.1)--(3.5)]{pichlershapiro}.

Policy consistency and equality of risk values are distinct as well. Xin and
Goldberg give moment-constrained inventory examples in which an initially
optimal policy ceases to be optimal after reoptimization, and another in which
the optimal policy sets agree although the two optimal values differ
\cite[Definitions~3.1--3.3, Observations~5--7, proofs in Sections~6.3--6.5]{xingoldberg}.
The company model must specify whether management can commit to a policy,
what information each later decision uses, and which family of joint laws is
retained. For a finite reference problem, enumerate complete admissible laws
and observation-measurable policies before choosing an approximation. For
larger problems, a set of beliefs over persistent types or an explicit
precommitment formulation can retain those restrictions, but its recursive
sufficiency and computational cost need a separate derivation. Repeated
worst-case optimization remains useful as a stress test when the changed
family is stated explicitly.

\subsection{Sequential Monte Carlo is an approximation to a stated model}
A particle approximation can reduce computation after a probability model is
chosen. It does not create that model. A type with probability $0.01$ has
probability $0.99^{100}>1/3$ of disappearing from 100 independent multinomial
resampling draws. No new evidence is involved in this loss. Rare but material
types therefore need explicit support checks, stratification or separate
conditional computations, with the resulting approximation stated.

For state-space models with persistent unknown parameters, SMC$^2$ combines
parameter particles and state filters, with particle-MCMC rejuvenation
\cite[Section~3.2, equation~(7), Proposition~1, Appendix~A]{smc2}. Its
extended target has the desired parameter posterior as a marginal under the
source assumptions. Finite particles still approximate that posterior; exactness
of a target marginal does not mean that one finite run represents every
possible type. The unbiased full particle likelihood estimate used in this
construction must be distinguished from an arbitrary estimated likelihood
ratio.

The inspected author-linked modern implementation creates a state filter for
each parameter particle, runs that filter through the current observation date
when evaluating a proposed parameter, and incorporates a likelihood correction
when increasing the state-particle count. These operations have real time and
memory costs. Retaining only filtering arrays differs from retaining complete
genealogies. This is a candidate for explicitly probabilistic filtering, not
a replacement for symbolic unweighted historical ambiguity.


\subsection{A finite filter that preserves contract identity}
Let $\theta$ be an unknown but fixed contract type, $x_t$ a changing hidden
operating state, and $y_t$ the analyst's observation. A declared probabilistic
model supplies transition probabilities $P_\theta(x'\mid x)$ and observation
likelihoods $\ell_t(y\mid\theta,x')$. With joint posterior weights
$w_{t-1}(\theta,x)$, the prediction and update are
\begin{align}
w_t^-(\theta,x')&=\sum_x
  P_\theta(x'\mid x)w_{t-1}(\theta,x),\\
Z_t&=\sum_{\theta,x'}\ell_t(y_t\mid\theta,x')w_t^-(\theta,x'),\\
w_t(\theta,x')&=\ell_t(y_t\mid\theta,x')w_t^-(\theta,x')/Z_t.
\end{align}
There is no summation that changes $\theta$ into another contract type. A
contract modification needs an explicit event and a different transition
model. If $Z_t=0$, the result is an evidence conflict under the declared model;
normalizing arbitrary fallback weights would answer another question.

An uninformative observation has constant likelihood across states, so its
update does not change the predicted relative weights. A missing disclosure
is uninformative only when the observation model makes it so. If no prior is
justified, propagate a feasible set of joint $(\theta,x)$ pairs through the
transition relation and evidence restrictions. That procedure produces support,
not posterior probabilities. A declared set of priors can instead be updated
member by member; zero-likelihood members need an explicit conditioning rule,
and one must not replace the result by independently chosen period marginals.

A source correction can remove an earlier restriction. Retaining only the
last posterior may then be insufficient: a zeroed type cannot be recovered
by multiplying it by another likelihood. Keep the original prior or feasible
family, the evidence versions and the transition model, and replay all
currently active evidence. Checkpoints can accelerate replay only when their
evidence dependencies remain valid. The reference filter compares this
recursion with direct path enumeration for four priors and all 27 length-three
observation sequences. It also checks rare positive types, uninformative
observations, conflict, and a correction that reopens a discarded type. This
is an exact finite method; estimating a defensible issuer likelihood and
controlling particle approximation remain separate problems.

\section{Financing must connect cash, promises and choices}
A financing transaction has three different descriptions. A contract creates
rights and promises; a market or bargaining model determines whether and on
what terms it is available; accounting rules determine recognition and later
measurement. A posting that debits cash and credits debt answers only the
last description for an assumed executed transaction. It does not establish
that lenders would advance the required cash.

\subsection{A consistent one-period debt convention}
Let $F$ be the promised payment at date $t+1$ and $B_t$ the cash advanced at
date $t$. The gross risk-free return is $R_f>0$. In this illustrative model
lenders are risk neutral and competitive under the specified probability
measure $Q$, the solvent probability is $1-d$, and default recovery paid at
$t+1$ is $R_D$. With no additional fees,
\begin{equation}\label{eq:debt}
B_t=\frac{(1-d)F+dR_D}{R_f},\qquad
q_t=\frac{B_t}{F},\qquad 1+y_t=\frac{F}{B_t}.
\end{equation}
The rate $y_t$ is a derived promised yield when $F$ and $B_t$ are positive.
Recovery and solvent payment are mutually exclusive states. Recovery paid
later requires discounting to its actual payment date. With a stochastic
discount factor under a physical measure, the corresponding general relation
is $B_t=\E_t[M_{t,t+1}\{(1-D)F+DR_D\}]$ under the stated pricing
assumptions. Physical default frequencies cannot be substituted for
risk-adjusted pricing probabilities without justification.

Take $F=100$, $d=0.1$, $R_D=40$ and $R_f=1.05$. Expected repayment is 94,
so $B_t=1880/21\simeq89.524$, $q_t\simeq0.89524$ and
$y_t=11/94\simeq11.702\%$. These numbers do not provide 95 of cash,
even though the promised repayment is 100. The simulator must return a
funding shortfall, search another admitted issue, reduce expenditure or invoke
a stated distress policy. Raising face value may change default probability
and recovery; it cannot be treated as a costless balancing operation.

A principal-based notation is also valid. If $b$ denotes cash advanced,
$F=b(1+y)$ and lender break-even reads
\begin{equation}
bR_f=(1-d)b(1+y)+dR_D.
\end{equation}
The problem in the supplied risky-debt note is the simultaneous use of $b$
as face in $b/(1+y)$ and as principal in this break-even equation. Those
statements do not define the same quantity. In the example, solving the
principal equation with $b=100$ and then treating that same 100 as face
gives cash $9000/101$. Its risk-free accumulation differs from 94 by
$-44/101$. The mismatch is wrong relative to lender break-even. Either
convention is repairable by defining face, proceeds and recovery consistently.
The local correction does not establish an erratum in the source edition of
Strebulaev and Whited \cite{strebulaev}; that version question remains open.

\subsection{Prices depend on the policy that follows issuance}
For longer-lived debt, lender receipts depend on future issuance, investment,
priority changes and default choices. A candidate policy $\pi$ produces a
price $q^\pi$. If management optimizes while holding an incompatible price
fixed, the resulting policy may not be the one lenders priced. A simultaneous
solution must make the chosen policy optimal under its price and make the
price value that policy's promised and default-state payments. A constrained
contract-price input is a useful alternative for mechanical stress testing,
but its interpretation is conditional on that price and availability.

Benzoni and coauthors make this distinction explicit in a dynamic model with
fixed issuance costs \cite[Section~2, Propositions~1--2, Appendix~A]{benzoni}.
Their debt quantity is face value, its price is per unit of face, and
shareholders choose default, issuance or inaction. Value matching and smooth
pasting generate a candidate barrier policy; off-policy deviations and a
supermartingale condition are required for equilibrium verification. The
retained model assumes zero default recovery and risk-neutral agents. It
therefore illustrates the policy-price consistency problem, but is not the
positive-recovery solution required by the preceding example.

The alternative managerial model of Lambrecht and Myers places investment,
borrowing, payout and rents in a dated sources-and-uses constraint and imposes
a no-Ponzi restriction \cite[Sections~2--4]{lambrecht}. Its safe-debt
solutions depend strongly on managerial preferences and governance. In its
exponential-utility example, the liquidation covenant protects lenders under
the stated diffusion dynamics. A jump across the boundary, delayed liquidation
or costly asset disposal would require a different recovery analysis. A
closed-form policy derived under those restrictions is a competitor to
consider, not an identified default policy for every issuer.

\subsection{Recovery, conversion and creditor disagreement}
Suppose a liquidation estate has value $A\geq0$, senior promised payment $S$
and junior promised payment $J$. Under an expressly assumed strict priority
contract with no liquidation costs or other claims,
\begin{equation}
R_S=\min(A,S),\quad
R_J=\min((A-R_S)^+,J),\quad
R_E=(A-R_S-R_J)^+.
\end{equation}
At $A=70$, $S=60$, $J=50$, creditors receive 60 and 10 and equity receives
zero. Applying $\min(A,J)$ independently to the junior claim would allocate
the same estate twice. This waterfall is a contractual illustration; actual
priority, collateral, netting and insolvency costs require the applicable
facts and authority. Consolidation does not make assets of every subsidiary
available to every creditor.

Conversion changes rights before a subsequent payoff is allocated. If junior
debt of face 50 converts into 40\% of residual equity, and the remaining
enterprise value is 100 subject to senior debt of 60, equity is worth 40.
The converted holder receives value 16 and the prior equity receives 24.
The junior debt has been extinguished. Paying the old 50 promise as well as
the new shares would double count the conversion. A covenant trigger, an
exercise decision and the register update must be dated separately; a
contract may specify a different conversion ratio or priority.

Recovery conservation does not settle bargaining. Consider an explicitly
authored restructuring example: liquidation yields 50 to a secured creditor
and zero to an unsecured creditor; continuation is worth 70 in total. An
offer of 45 to the secured creditor and 25 to the unsecured creditor has
higher total value than liquidation but will be rejected by a secured creditor
that can obtain 50 by refusing. Whether the offer can proceed depends on
voting, enforcement, bargaining and financing rules, which this arithmetic
does not supply. Thus a simulator must not choose the maximum-total-value
allocation and call it an implementable restructuring.

Dewatripont, Legros and Matthews provide a primary contracting comparison
\cite[Sections~2--6 and Appendix~B]{dewat}. Their environment has observable
but noncontractible effort, limited liability, monotone investor receipts and
efficient renegotiation. Under entrepreneur-offer renegotiation, a debt
contract can dominate a more general contract in the stated pure-strategy
class. Their appendix also supplies cases in which changed bargaining power,
third parties or mixed messages change that conclusion. Even their investor
conversion option need not be an ordinary equity security, as footnote 16
explains. Contract labels cannot replace the exact contingent payments and
exercise rights in the company model.

\subsection{New equity is a transfer of ownership as well as cash}
Let the pre-issue equity value be $V$, gross new subscription be $I$ and an
issue cost paid by the company be $f$. Assume homogeneous shares, immediate
execution, no other valuation effect and fairly priced new shares. Post-issue
equity value is $V+I-f$. If new investors receive fraction $\alpha$, fairness
requires
\begin{equation}
\alpha(V+I-f)=I,\qquad
\alpha=\frac{I}{V+I-f}.
\end{equation}
Old investors retain value $(1-\alpha)(V+I-f)=V-f$. For $V=80$, $I=20$
and $f=2$, the new fraction is $10/49$ and old equity is worth 78. With
100 old shares, a common post-issue price is $0.78$ and new investors receive
$1000/39$ shares. Their value is 20 and total equity is 98.

This result places the issue cost on incumbent shareholders under the stated
fair-pricing convention. It is not the universal price convention for an
offer announced before costs are known. If $f\geq V$, a positive common
share price and this issuance convention cannot be maintained. Preferential
rights, discounts, control premiums, warrants, convertibles, restrictions on
fractional shares and bargaining need explicit terms. Book equity, market
equity value and share count must remain distinct.

\subsection{A rate shock changes both market access and ownership}
An exact combined example shows why the financing steps must be dated. Two
possible companies each start with cash 16 and loan face 100. One pays a
fixed coupon of 5; the other pays 7 after the specified floating-rate reset.
Each then makes an operating payment of 6. A covenant tests cash of at least
5 immediately after these payments and before financing. Breach forbids new
borrowing but permits an equity cure. Management requires closing cash 8.
These are assumed contract terms, not a claim about a standard covenant.

The fixed-rate company has cash $16-5-6=5$, meets the covenant and needs 3.
A candidate debt issue with maximum face 5, conditional solvency probability
$4/5$, recovery 2 and gross risk-free return $21/20$ has proceeds $88/21$.
Holding that quoted per-face price fixed for this conditional stress, issuing
face $315/88$ raises exactly 3. The floating-rate company has cash
$16-7-6=3$, breaches the covenant and needs 5. Debt is unavailable; even an
improper borrowing override at the maximum face would raise only $88/21<5$.

Suppose the permitted equity cure has fee 1, pre-issue economic equity value
20 and 100 incumbent shares. A gross subscription of 6 provides net cash 5.
The post-issue equity value is 25, new investors own $6/25$, and the common
share price is $19/100$. Their $600/19$ new shares have value 6, while
incumbents retain value 19. Both companies finish with cash 8, but their new
debt and ownership differ. Merging them on closing cash would conceal those
differences. Testing the covenant only after the cure would erase the earlier
breach and wrongly admit the borrowing alternative.

This example holds quoted financing prices and the pre-issue equity valuation
fixed to expose event ordering and allocation. It does not solve the feedback
from financing to default and prices. A source correction introduces another
interaction: withdrawing an unsupported fixed-rate classification can reopen
the floating-rate history, invalidating an earlier single-history result.
If management later receives a rate notice, its admissible policies can change
even when the analyst's reported balances do not. Shared computations must
therefore depend on evidence version, actor information, contracts and the
requested outcomes, as well as current cash.



\subsection{Executing the distinction between promises and book amounts}
The combined example can be executed through the existing typed company
runtime once its state distinguishes contractual face from carrying amount.
For the fixed-rate history, the authored initial-measurement convention posts
cash of 3 and a liability carrying amount of 3, while contractual face increases
by $315/88$. The closing liability carrying amount is 103 and outstanding face
is $9115/88$. Subsequent effective-interest measurement and settlement require
further methods; substituting carrying amount for the promise would lose the
very exposure the example is intended to preserve.

For the floating-rate history, subscription proceeds of 6 increase cash and
contributed capital; the direct issue cost of 1 reduces both. The typed state
retains 100 incumbent shares and $600/19$ new shares. Under this explicitly
synthetic posting convention, both histories report assets of 112. The fixed
history has liabilities of 103 and book equity of 9; the floating history has
liabilities of 100 and book equity of 12. These book amounts differ from the
economic equity values used to price the subscription. No reporting-standard
admission follows from the balanced journals.

The reference adapter executes coupon payment, operations, covenant observation
and financing in that order, using the unchanged common company executor.
It rejects borrowing after the floating-rate breach, proceeds inconsistent
with the supplied per-face quote, a zero issue price, duplicate issuance and
payment without sufficient cash. A mutation that omits the share-count update
still produces balanced accounts and completes the accounting path; the
independent ownership comparison detects the error. Accounting conservation
alone therefore cannot validate a financing method.

The same runtime also executes the earlier information example. Starting with
cash 6, a fixed coupon 5 leaves 1, while a floating coupon 7 cannot be paid.
A hedge costing 1 caps the coupon at 5. An informed manager hedges only the
floating history and retains cash 1 or 0; a manager who must act identically in
both histories can hedge both and retain 0. Withdrawing the earlier fixed-only
evidence reopens the payment-refusal branch. The negative cash in the algebraic
shortfall calculation denotes unmet funding; the runtime refuses the payment
atomically and keeps the opening cash, rather than recording an impossible
settlement. A missed-payment event and any ensuing default process would need
separate declared methods.

\subsection{Checking lender prices against the issuer's actual choice}
In a finite one-period model, policy consistency can be checked by enumeration.
Specify a candidate policy $\pi$, the distribution of the terminal estate under
that policy, the priority of each claim, and a pricing measure and discount
factor. Allocate the estate once, calculate each lender's expected discounted
receipt, and use those receipts as candidate issue proceeds. At those fixed
quoted proceeds, compare the issuer's payoffs from every feasible policy.
Retain the candidate only when its conjectured policy is a best response and
all funding and participation constraints hold. Ties must be retained. If no
pure-policy pair passes, report that result for the declared class; neither a
balancing loan nor an assumed mixed equilibrium repairs it automatically.

For a concrete local example, existing senior debt has face 40 and newly issued
junior debt has face 80. A project costs 60. The safe policy produces estate 100;
the risky policy produces estate 20 or 220 with equal probabilities. Both pricing
and shareholder valuation use the stated risk-neutral probabilities and gross
discount factor $21/20$. Absolute priority in this illustration gives receipts
\begin{equation}
(s,j,e)=\bigl(\min(A,40),\min((A-40)^+,80),(A-120)^+\bigr).
\end{equation}
Thus the safe estate pays $(40,60,0)$, and the risky estates pay $(20,0,0)$
and $(40,80,100)$. Each allocation sums to its estate. The candidate junior
proceeds are $400/7$ under the safe policy and $800/21$ under the risky policy.

At a posted quote $P$, existing shareholders must subscribe $60-P$ to cover
the investment. Their net value from policy $\pi$ is
$\E[e_\pi]/(21/20)-(60-P)$. They prefer risk at either quoted price.
Consequently, the safe-policy quote is inconsistent with their best response:
if risk is chosen, junior lenders receive present value $800/21$ after advancing
$400/7$, a shortfall of $400/21$. The risky quote is consistent within the
pure-policy class, but requires shareholder cash $460/21$. If shareholders
have only 20, the project remains unfunded. The safe quote would need less cash,
but its incentive problem remains. Even an enforceable safe-policy commitment
would still need shareholder participation: here its zero terminal equity
receipt does not cover their positive subscription. Liquidity, incentives,
lender valuation and voluntary participation are distinct restrictions.

The same separation applies to restructuring. Given estate 70 and creditor
outside values 50 and 20, an offer 45/25 conserves value but the senior creditor
refuses. An offer 50/20 meets both constraints. A mandatory fee 1 reduces the
available estate to69 and makes simultaneous participation impossible under
these stated outside values. This is a feasibility result. Selecting a unique
offer requires bargaining rights or another decision rule; it does not follow
from total value conservation.

\subsection{Positive continuation value need not make renegotiation feasible}
Hart and Moore's retained 1991 working paper makes the distinction between
productive value and credible repayment particularly clear. The entrepreneur
can withdraw human capital, the creditor can seize specified assets, and
renegotiation follows an explicit offer game. Returns and liquidation values
are deterministic, private savings have a particular seizure treatment, and
one creditor controls the assets \cite[Sections~2--3 and Appendix~A1--A4]{hartmoore}.
Under those assumptions, Lemma~1 derives the creditor's continuation reservation
value and requires both future nonrepudiation and cash feasibility. The bargaining
rule determines a result that cannot be obtained merely by balancing accounts.

In their Example~1, take two remaining operating receipts of 2 and 26 and
liquidation rents of 8 in each period. The creditor's reservation value is
16, while the final credible payment is at most 13. The entrepreneur therefore
needs to provide at least 3 before the final date. The first operating receipt
provides only 2, so private wealth $w\ge1$ is necessary. At $w=1$, an immediate
payment of 1 followed by payments 2 and 13 satisfies the cash and repayment
conditions. At $w<1$, it cannot. Liquidation occurs despite operating value
28 exceeding liquidation value 16 \cite[Example~1, printed p.~16]{hartmoore}.
The reference calculation reproduces this boundary with exact rational amounts.
It is a useful counterexample to assuming that positive total surplus guarantees
refinancing or restructuring.

Their repayment-path analysis also reinforces the distinction between an
observed payment and its contractual explanation. A long-term contract and
rolling short contracts can support the same path under the paper's conditions;
the latter may additionally require access to outside refinancing
\cite[Proposition~2 and printed pp.~27--28]{hartmoore}. A future market closure
can therefore distinguish histories that looked equivalent before the shock.
For a multi-creditor company, the analyst must additionally specify priority,
control and consent rights, collateral and intercreditor restrictions. The
single-creditor result supplies neither an equity-issuance model nor a general
legal restructuring procedure.

\section{What is disclosed is part of the observation model}
Imagine equal prior probabilities of a good and bad type. If good types
disclose with probability $0.9$ and bad types with probability $0.1$, then
disclosure raises the probability of a good type to $0.9$, while silence
reduces it to $0.1$. Treating a missing disclosure as no information would
retain $0.5$ in both cases. The discrepancy comes from omitting the selection
mechanism, not from numerical approximation. These probabilities are an
authored example, not estimates of actual issuer disclosure behaviour.

Let $Y$ denote potentially observable information, $M$ the pattern of
observed components, $f_\theta(Y)$ its model and $g_\phi(M\mid Y)$ the
publication or missingness mechanism. The likelihood of the realised data is
\begin{equation}
L(\theta,\phi)=\int
f_\theta(y_{\rm obs},y_{\rm mis})
g_\phi(m\mid y_{\rm obs},y_{\rm mis})\,dy_{\rm mis}.
\end{equation}
Here the realised pattern $m$ fixes which coordinates are observed; the
integration is over the missing coordinates only. Discrete missing components
use a sum instead. If, for the realised observation, $g_\phi(m\mid y)$ is constant over
all compatible missing components, it factors out of this integral. Together
with distinct parameter spaces this yields the relevant likelihood
proportionality; a Bayesian factorisation also needs the appropriate prior
independence. Seaman and coauthors distinguish these realised-MAR conditions
from the everywhere-MAR conditions relevant to repeated-sampling comparisons
\cite[Theorems~1--3, Section~5.4 and Appendix]{seaman}. These are sufficient
conditions for specified inferential equivalences. They do not demonstrate
that corporate missingness satisfies them.

\subsection{A simple equilibrium of silence}
A separate local illustration makes non-disclosure endogenous. Let value $V$
be uniform on $[0,1]$. A manager knows $V$ with probability $1/2$, independently
of value; otherwise the manager cannot disclose it. Informed disclosure is
truthful and costless, the manager maximizes the contemporaneous price, and
risk-neutral investors price at conditional expected value. Suppose an informed
manager discloses above cutoff $a$. Silence occurs with probability
$(1+a)/2$. Its unconditional value contribution is
$1/4+a^2/4$, giving silence price
\begin{equation}
P_{\rm silence}(a)=\frac{1+a^2}{2(1+a)}.
\end{equation}
An informed manager at the cutoff is indifferent, so $a=P_{\rm silence}(a)$.
Multiplying through and collecting terms gives $a^2+2a-1=0$, hence
$a=\sqrt2-1\simeq0.4142$. Types above that value disclose and those below
it prefer the silence price. The uninformed manager explains why silence
does not identify the worst type. This derivation is a local example; the
blocked Dye original is not claimed to have been reproduced.

Nor is favorable selection universal. Suijs models a firm that knows its
return but is uncertain how investors compare it with an outside opportunity
\cite[Sections~3--4, Proposition~2 and Appendices~A--B]{suijs}. Under the
stated conditions a profitable firm can disclose lower returns and withhold
higher returns. Multiple partial-disclosure equilibria are possible. The
implication for AccDSL is to represent alternative disclosure mechanisms and
test their consequences, rather than hard-code one direction of selection.

\subsection{Mechanical shocks, conditioning and intervention}
A mechanical rate stress replaces an input while holding specified other
relations fixed. A conditional forecast computes outcomes among states where
an observed rate has a chosen value. A causal intervention changes the
mechanism generating the rate and then propagates the remaining structural
relations. These calculations can disagree.

For a finite counterexample, let $U$ be a fair binary variable and observe
$X=U$, $Y=U$. An alternative structural model has $X=U$, $Y=X$. Both give
exactly the same observed joint distribution. Under an intervention setting
$X=1$, the first model retains $\E[Y]=1/2$ while the second gives
$\E[Y]=1$. Observational agreement does not identify the intervention
response. A forecast labelled as causal therefore needs the structural
assumptions that distinguish these possibilities, or must retain the set of
responses. Source corrections and simulated observations also have different
time meanings: a hypothetical future disclosure updates an actor inside that
run, not the analyst's actual present evidence.

\section{Contract algorithms and learned solvers have narrower jobs}
\subsection{What a contract engine establishes}
A contract engine can generate scheduled events and payoffs from supplied
terms, calendars and risk observations. It need not determine which financing
is available, choose management's action, infer undisclosed terms or price
endogenous default. The retained ACTUS specification identifies itself as
v1.1, revision 843f7a3, dated 8 June 2020, aligned with dictionary v1.2
\cite{actus}. That historical pairing is more precise than treating a current
website download as the latest compatible runtime.

Independent arithmetic on the retained official PAM examples checks two
simple schedules. A principal of 3000 at 10\% over 31 actual days with A365
pays $3000(0.1)(31/365)=25.479452\ldots$. The discounted-issue A360 case
advances 2800 against face 3000 and pays $3000(0.1)(59/360)$ on its first
two-month interest date. Initial interest of zero and interest-before-maturity
ordering are retained from the explicit schedules. These calculations match
24 listed events within the stated decimal-serialization tolerance.

This does not execute ACTUS. The inspected service requires ACTUS core version 1.1.0;
its Dockerfile copies and builds a local
core directory. The attempted public core repository and Maven lookup did
not yield the matching implementation. A historically matched dictionary and
specification therefore remain insufficient to claim runtime conformance.
Rate resets, business-day adjustment, missing observations and other contract
types need their own implementation traces and comparisons before adoption.

\subsection{The equation being minimized matters}
Suppose an economic equation is $\E[r(\theta,\epsilon)\mid s]=0$. The
squared conditional mean and conditional mean square are different:
\begin{equation}
\E[r^2\mid s]=\{\E[r\mid s]\}^2+\operatorname{Var}(r\mid s).
\end{equation}
The variance can depend on the policy parameters. In the exact example
$r=\theta-1+\theta\epsilon$, with independent fair $\epsilon\in\{-1,1\}$,
the intended loss is $(\theta-1)^2$, minimized at 1. The mean square is
$(\theta-1)^2+\theta^2$, minimized at $1/2$. Training can successfully
optimize the wrong quantity.

Conditionally independent draws $\epsilon_a,\epsilon_b$ give
\begin{equation}
\E[r(\theta,\epsilon_a)r(\theta,\epsilon_b)\mid s]
=\{\E[r(\theta,\epsilon)\mid s]\}^2.
\end{equation}
Conditional independence, the conditioning state and differentiation
assumptions matter. Reusing the same draw returns the mean square. A finite
product estimate can be negative even though its expectation is nonnegative.
Maliar, Maliar and Winant discuss solving economic equations with deep
learning \cite{maliar}; the inspected author-linked Krusell--Smith notebook
uses the product of residuals from two future draws. That code is evidence
for its equation construction, not evidence that it solves a risky-debt
issuance and recovery system.

Gradients require the same care. If $F(\theta,q)=\theta q$ and
$q(\theta)=\theta$, the total derivative is $2\theta$. Differentiating
with $q$ detached gives $\theta$. A deliberate alternating-price algorithm
may use the latter as one step, but it is wrong if described as the total
gradient of the coupled objective. Boundary default decisions and discrete
issuance choices add feasibility and nonsmoothness issues beyond this example.

The notebook also contains differences between explanatory markdown and
executed parameter values and initialization, along with clipping and
cross-sectional normalization of shocks. Those change the problem being
computed and must be included in a reproduction. The separate audit retains
exact cells and hashes. No neural training has been promoted here. The
Duarte paper's announced replication library has not been matched to an
author-linked implementation in the inspected sources \cite{duarte}.


\subsection{What a causal generator does and does not preserve}
Acciaio, Eckstein and Hou construct a time-causal variational autoencoder whose
encoder and decoder at date $t$ use only input coordinates available by $t$.
Their retained preprint bounds a \emph{directional} causal transport discrepancy
by reconstruction error and a term involving latent relative entropy
\cite[Lemma~2.4, Theorem~2.5 and Corollary~2.6]{tcvae}. The relative-entropy
step requires the reconstructed and generated distributions to have the stated
bounded support. Normalizing the observed sample does not establish that a
learned decoder satisfies that support restriction for every generated input.
A finite training estimate is also different from a population bound.

There is a further issue in transferring the control result. Theorem~3.1 of
the retained version states its comparison for a closed convex subset
$\mathcal H$ of adapted strategies. Its proof takes a policy $G$ on the second
process and constructs $\widetilde G(x)=\E[G(Y)\mid X=x]$. Causality establishes
adaptedness, but adaptedness alone does not establish
$\widetilde G\in\mathcal H$. Here is a direct counterexample to that stated
generality. There is one date, $\mu=\delta_1$, $\nu=\delta_0$, and
\begin{equation}
\mathcal H=\{h\},\qquad h(x)=\min(1,\max(0,x)),\qquad Q(x,a)=a.
\end{equation}
The singleton class is closed and convex, and $Q$ is convex in $a$ and
uniformly zero-Lipschitz in $x$. Nevertheless,
$V(\mu)-V(\nu)=1$, whereas the proposed right-hand side with $L=0$ is zero.
The transferred action at $x=1$ is zero, outside the singleton policy class.
This disproves the unrestricted policy-class statement in that retained
version. It does not refute the argument when the feasible class is explicitly
closed under the conditional transfer. The published edition has not been
matched equation by equation, so this is not asserted to be a journal erratum.

For AccDSL, state-dependent borrowing limits, integer issuance sizes and
contractual exercise restrictions make closure substantive. A repaired transfer
argument must verify feasibility of the mapped policy, or use a method that
retains the original feasible actions. The one-sided value comparison is not
itself a bound on the performance of the same learned policy after deployment.
The inspected paper-linked code has forward LSTM encoder/decoder paths, but
also external conditioning variables; those variables require their own
availability dates. Its selected conditional decoder uses ReLU and a final
linear layer, so the architecture does not enforce the theorem's unit-ball
support. These observations identify additional obligations before adoption;
no training result is inferred from source-code inspection.

\subsection{Learning a transition law needs more than matching an aggregate}
Sald\i{} and Y\"uksel distinguish weak and strong kernel-embedding topologies
and give a square-wave counterexample to their equivalence
\cite[Theorem~4.1 and proof]{kerneltop}. Let $Y$ be uniform on $[0,1]$.
On alternating half-cells of a fine partition, make a binary successor $U$
deterministically 1 or 0. As the partition becomes finer, smooth integrated
tests can approach those of a fair binary kernel. At every particular input,
however, the conditional law remains a point mass. For the Kronecker kernel
on $\{0,1\}$, its squared embedding discrepancy from the fair law is $1/2$.
The local exact calculation checks this discrepancy and the vanishing error
for a linear test, rather than claiming that a few tests prove weak convergence.

Their discounted-control application uses bounded continuous costs, finite
actions, continuous transition kernels and a common dominating measure, including
for the initial distribution. One route requires a strong essential-supremum
condition; the weaker route adds bounded equicontinuous densities
\cite[Assumptions~2--4, equation~(11), Theorems~5.1--5.2]{kerneltop}.
A default discontinuity, hidden creditor information or an action-dependent
financing constraint therefore needs a separate argument. An equally important
problem is support: models that agree under the historically chosen action may
disagree under a new action. With action losses $(1,100)$ in the real model and
$(1,0)$ in an approximate model, historical use of the first action gives
identical observations. Optimizing the approximation chooses the second action
and incurs regret 99. Agreement under the historical policy alone cannot
establish a policy-uniform approximation claim.


The original robustness result used in that application makes another condition
explicit. Kara and Y\"uksel distinguish continuity of optimal values from
robustness of the policies attaining them. Their discounted-control argument
uses continuous kernel convergence, bounded continuous costs and compact
actions. Policy transfer additionally requires each approximate policy to be
optimal from every initial state, as when it satisfies the discounted Bellman
equation throughout the state space
\cite[Assumption~1.1, Theorems~4.2--4.4 and Appendix~A.5]{kara}.
An optimizer trained only from a historical starting state can leave its actions
elsewhere arbitrary; the true transition can reach those neglected states.
The paper's counterexample makes precisely this distinction. Its finite-state
learning result also assumes continuing coverage of every state--action pair
\cite[Corollary~5.1]{kara}. Financial statements generated by one management
policy do not by themselves provide that coverage.

\section{What would establish an estimated company policy?}
Historical fit is only one restriction on a model. If two policies take the
same action on every observed state but diverge after refinancing is refused,
the historical sample cannot distinguish their response to that refusal.
The observation model can create another equivalence: infrequent disclosure
of losses can resemble a lower loss frequency. A source-to-parameter map
should therefore identify what is measured, what is inferred under a model,
and what remains set-identified or a sensitivity assumption.

A credible evaluation uses information as it was available at each forecast
origin. Later restatements and later-known covenant breaches remain unavailable
to earlier forecasts. Issuer exit and default must be retained rather than
dropping failed companies from the comparison. The benchmark ladder should
include a naive persistence/accounting baseline, a tuned classical policy or
state-space model, the plain proposed model and any enhanced version, using
the same information, feasible actions and evaluation periods.

The primary criterion must match the proposed use. For a probabilistic
forecast of default, the Brier loss $(q-D)^2$ has expectation
$p(1-q)^2+(1-p)q^2=p(1-p)+(q-p)^2$, uniquely minimized by $q=p$. This
identity is useful but does not validate the estimated $p$. Proper scoring
rules assess distributional forecasts under specified outcomes
\cite{gneiting}; unweighted sets require coverage and sharpness questions
appropriate to sets, not arbitrary probabilities added to permit scoring.
Decisions additionally need realized cost or regret evaluation under a
defensible comparison.

Repeated forecasts from one issuer and common macroeconomic shocks are
dependent. Paired comparisons, issuer/time blocking and an uncertainty method
appropriate to that dependence must be specified before ranking methods.
Rare defaults may make precise tail comparisons impossible in the available
panel. Such results remain descriptive. Conservation and feasibility errors
can veto a candidate regardless of its predictive score; a good score does
not repair a wrong accounting or financing quantity.


\section{A finite compression method that keeps discrete decisions}
A stronger but directly checkable condition can avoid averaging actions.
Consider a finite horizon of $T$ transitions. Keep a persistent contract type
$\theta$ fixed and let a history record the actor's observations and past
actions. In both the full model and a reduced model, require the same feasible
actions at paired histories. Couple their initial states so that the required
observations agree. At each paired history and common action, suppose a joint
coupling of the next state, observation and reported output leaves the states
paired and the observation/output records equal with conditional probability
at least $1-\delta_t$. This condition includes every history and action reachable
before a mismatch; fitting one historical policy is insufficient.

If $E_t$ is the event that no mismatch has occurred through transition $t$,
then $\Prob(E_t\mid E_{t-1})\ge1-\delta_t$. Repeated conditioning gives
\begin{equation}
\Prob(E_T^c)\le1-\prod_{t=1}^T(1-\delta_t)
\le\sum_{t=1}^T\delta_t.
\label{eq:finite-coupling}
\end{equation}
A common deterministic policy chooses the same action while the observation
histories agree. Randomized policies can use the same independent random draws,
if both models permit that randomization. No convexity or fractional issuance
is required. A bounded path loss $L\in[0,B]$ agrees on $E_T$, and hence
\begin{equation}
|\E L-\widetilde\E L|
\le B\Prob(E_T^c)
\le B\left\{1-\prod_{t=1}^T(1-\delta_t)\right\}.
\label{eq:finite-cost}
\end{equation}
A default indicator is covered because it is a bounded output and the
coupling condition compares its probability explicitly. Small cash distance
alone does not provide that condition. If exact output matching is impossible,
a separate bound on the loss difference on matched paths must be added.

For finite kernels with common outcome labels, the smallest one-step mismatch
probability is their total-variation distance
$\frac12\sum_z|K(z)-\widetilde K(z)|$. Pairing full states with reduced states
requires a coupling whose successful pairs remain in the declared relation;
marginal observation agreement is insufficient. An implementable finite
procedure first partitions states by persistent type, available information,
feasible actions and required outputs. It then compares joint successor
kernels into the current blocks for every action, splitting a block whenever
these differ. Repeat until no block changes. Each split reduces the number of
states per block, so a finite partition stabilizes after finitely many splits.
This yields an exact quotient for the declared finite model, not a claim that
the model contains every economically possible company.

Approximate sharing uses the same comparison with an explicit error allowance
and a checked pairing. Threshold regimes and different feasible actions are
kept separate unless their mismatch has been bounded. An unresolved comparison
is retained as unresolved, not assigned zero error. For a fixed common policy
class, the uniform cost bound in \eqref{eq:finite-cost} also gives the earlier
$2\varepsilon+\eta$ policy-regret bound. If the persistent types are unweighted,
apply the bound within each type. A uniform type-wise bound also bounds the
worst-type value, since
$|\sup_\theta J_\theta-\sup_\theta\widetilde J_\theta|
\le\sup_\theta|J_\theta-\widetilde J_\theta|$.
The adversary does not redraw the contract type at each date.

The exact reference uses three periods, cash observed by management, a hidden
fixed type, binary borrowing, and a default penalty. It compares complete
paths for all 64 policies in the explicitly restricted class of date-and-
zero-cash rules. Removing an irrelevant tag preserves every path. Merging
types is rejected. A synthetic change of $1/20$ in each period's bad-expense
probability gives the conservative cost bound
$12\{1-(19/20)^3\}=3423/2000$. The largest observed value discrepancies are
$4433/8000$ and $1/2$ in the two types, within the bound. These finite checks
support the stated method and expose its conservatism. They do not supply a
capacity estimate or a bound for an unenumerated continuous company model.

\section{A method with explicit limits}
The resulting design retains joint economic histories, dated evidence and
actor-specific information. It uses symbolic constraints and shared records
before approximation; exact merging requires equality of the relevant
continuation law. Where approximation is needed, cash, threshold events and
policy regret receive different error arguments. Probabilistic filters operate
only after prior, transition and observation models have been stated.
Financing connects proceeds to contingent creditor payments, management
decisions and ownership changes before generating the corresponding journal.

Several residuals are substantive. The historical planning sources still
need source-specific access and edition comparison. The full ACTUS runtime
has not been matched. General creditor bargaining, broad hybrid-state
approximation and empirical issuer-policy identification remain open.
The acquired follow-up literature also contains work awaiting technical
inspection. These limitations prevent an all-gap closure claim. They do not
erase the exact corrections, source-supported restrictions and counterexamples
that can already govern the next implementation.

\begingroup
\raggedright
\begin{thebibliography}{99}
\bibitem{mccarthy} W. E. McCarthy (1982). The REA Accounting Model: A Generalized Framework for Accounting Systems in a Shared Data Environment. \emph{The Accounting Review} 57(3), 554--578. Retained author-hosted scan.
\bibitem{pflug} G. C. Pflug and A. Pichler (2012). A Distance for Multistage Stochastic Optimization Models. \emph{SIAM Journal on Optimization} 22(1), 1--23. doi:10.1137/110825054.
\bibitem{adapted} J. Backhoff-Veraguas, D. Bartl, M. Beiglböck and M. Eder (2020). Adapted Wasserstein distances and stability in mathematical finance. \emph{Finance and Stochastics} 24, 601--632. doi:10.1007/s00780-020-00426-3.
\bibitem{estimation} J. Backhoff, D. Bartl, M. Beiglböck and J. Wiesel (2022). Estimating processes in adapted Wasserstein distance. \emph{Annals of Applied Probability} 32(1), 529--550. doi:10.1214/21-AAP1687. Technical anchors refer to retained arXiv:2002.07261.
\bibitem{epstein} L. G. Epstein and M. Schneider (2007). Learning under ambiguity. \emph{Review of Economic Studies} 74, 1275--1303. doi:10.1111/j.1467-937X.2007.00464.x. Retained technical reading in the preceding dossier.
\bibitem{smc2} N. Chopin, P. E. Jacob and O. Papaspiliopoulos (2013). SMC$^2$: an efficient algorithm for sequential analysis of state space models. \emph{Journal of the Royal Statistical Society B} 75, 397--426. doi:10.1111/j.1467-9868.2012.01046.x. Technical anchors refer to retained arXiv:1101.1528.
\bibitem{strebulaev} I. A. Strebulaev and T. M. Whited (2012). Dynamic Models and Structural Estimation in Corporate Finance. \emph{Foundations and Trends in Finance} 6(1--2), 1--163. Retained edition and local risky-debt convention comparison; correction correspondence unresolved.
\bibitem{benzoni} L. Benzoni, L. Garlappi, R. S. Goldstein and C. Ying (2023). Debt Dynamics with Fixed Issuance Costs. Federal Reserve Bank of Chicago Working Paper WP 2023-01. doi:10.21033/wp-2023-01. Retained file includes the online appendix by J. Hugonnier.
\bibitem{lambrecht} B. M. Lambrecht and S. C. Myers. The Dynamics of Investment, Payout and Debt. Retained Cambridge manuscript dated 26 January 2017; journal identity doi:10.1093/rfs/hhx081.
\bibitem{dewat} M. Dewatripont, P. Legros and S. A. Matthews (2003). Moral Hazard and Capital Structure Dynamics. PIER Working Paper 03-006, January 29 version. Journal identity doi:10.1162/154247603322493186; no edition-equivalence assertion.
\bibitem{seaman} S. Seaman, J. Galati, D. Jackson and J. Carlin (2013). What Is Meant by ``Missing at Random''? \emph{Statistical Science} 28(2), 257--268. doi:10.1214/13-STS415. Retained electronic reprint arXiv:1306.2812.
\bibitem{suijs} J. Suijs (2007). Voluntary disclosure of information when firms are uncertain of investor response. \emph{Journal of Accounting and Economics} 43, 391--410. doi:10.1016/j.jacceco.2006.10.002. Retained Tilburg repository copy.
\bibitem{actus} ACTUS Financial Research Foundation. Technical Specifications v1.1, revision 843f7a3, 8 June 2020; declared alignment with dictionary v1.2. Official specification, dictionary, fixture and service snapshots retained separately.
\bibitem{maliar} L. Maliar, S. Maliar and P. Winant (2021). Deep learning for solving dynamic economic models. \emph{Journal of Monetary Economics} 122, 76--101. doi:10.1016/j.jmoneco.2021.07.004. Author-linked Krusell--Smith notebook audited at commit 4f086d9a2586969406a00f88d3b036ba1413b8d3.
\bibitem{duarte} V. Duarte and J. Fonseca (2026). AI for Structural Estimation. NBER Working Paper 35283, May. Retained paper and author page. The announced replication package is not matched to inspected author-linked code; technical claims remain limited to the preceding paper reading.
\bibitem{gneiting} T. Gneiting and A. E. Raftery (2007). Strictly Proper Scoring Rules, Prediction, and Estimation. \emph{Journal of the American Statistical Association} 102, 359--378. doi:10.1198/016214506000001437. Retained technical reading in the future-validation dossier.
\bibitem{shapiro} A. Shapiro. Rectangular Sets of Probability Measures. Retained 24-page Optimization Online manuscript; journal identity: \emph{Operations Research} 64(2), 528--541 (2016), doi:10.1287/opre.2015.1466. Technical anchors refer to the retained manuscript; edition equivalence is not asserted.
\bibitem{pichlershapiro} A. Pichler and A. Shapiro (2021). Mathematical Foundations of Distributionally Robust Multistage Optimization. arXiv:2101.02498v2, 13 May 2021. Technical anchors refer to this retained version.
\bibitem{xingoldberg} L. Xin and D. A. Goldberg. Time (in)consistency of multistage distributionally robust inventory models with moment constraints. Retained 43-page Optimization Online manuscript. Journal identity: \emph{European Journal of Operational Research} (2021), doi:10.1016/j.ejor.2020.07.041. Technical anchors refer to the retained two-author version, not the earlier three-author preprint.

\bibitem{hartmoore} O. Hart and J. Moore (1991). A Theory of Debt Based on the Inalienability of Human Capital. NBER Working Paper 3906, November. doi:10.3386/w3906. All 62 pages of the retained scan inspected; not asserted equivalent to the 1994 journal edition.
\bibitem{probview} M. Beiglb\"ock, S. Pfl\"ugl and S. Schrott (2024). A Probabilistic View on the Adapted Wasserstein Distance. arXiv:2406.19810v1. Technical anchors refer to this version; journal identity doi:10.1016/j.spa.2026.105032.
\bibitem{tcvae} B. Acciaio, S. Eckstein and S. Hou (2024). Time-Causal VAE: Robust Financial Time Series Generator. arXiv:2411.02947v1. Journal identity doi:10.1137/24m1711650; version-specific transfer issue discussed in the text. Author-linked code inspected at commit a43effe0125d7c778d0560e4704754991d8e271b.
\bibitem{kerneltop} N. Sald\i{} and S. Y\"uksel (2025). Kernel Mean Embedding Topology: Weak and Strong Forms for Stochastic Kernels and Implications for Model Learning. arXiv:2502.13486v2. Journal identity doi:10.1214/25-aap2266; technical anchors refer to the preprint.
\bibitem{kara} A. D. Kara and S. Y\"uksel (2020). Robustness to Incorrect System Models in Stochastic Control. Retained arXiv:1803.06046v7, 1 February 2020. Journal identity: \emph{SIAM Journal on Control and Optimization} 58, 1144--1182, doi:10.1137/18M1208058. Technical anchors refer to the retained preprint.
\end{thebibliography}
\endgroup
\end{document}
